\chapter{CALCULOS JUSTIFICATIVOS}

\section{Calculo de la demanda maxima}

\subsection{Premisas de calculo}

Para el presente trabajo se adopta una metodologia simplificada por circuitos. La potencia instalada se estima a partir de cargas tipicas residenciales y la maxima demanda se obtiene aplicando factores de demanda preliminares acordes con el tipo de carga.

\subsection{Formula general}

\begin{equation}
P_{inst} = \sum P_i
\end{equation}

\begin{equation}
P_{dem} = \sum (P_i \times f_{d,i})
\end{equation}

\begin{equation}
I_{dem} = \frac{P_{dem}}{V \times fp}
\end{equation}

Donde:

\begin{itemize}
  \item $P_{inst}$: potencia instalada total.
  \item $P_{dem}$: potencia de demanda total.
  \item $P_i$: potencia instalada por circuito.
  \item $f_{d,i}$: factor de demanda del circuito.
  \item $I_{dem}$: corriente demandada.
  \item $V$: tension nominal.
  \item $fp$: factor de potencia.
\end{itemize}

\subsection{Cuadro de cargas por circuito}

\begin{table}[H]
\centering
\caption{Cuadro de cargas por circuito}
\label{tab:cuadro-cargas-circuito}
\begin{tabularx}{\textwidth}{p{1.6cm}p{4.2cm}r r r r X}
\toprule
\textbf{Circuito} & \textbf{Uso} & \textbf{Pot. inst. W} & \textbf{Factor} & \textbf{Pot. dem. W} & \textbf{Corriente A} & \textbf{Observacion} \\
\midrule
C1 & Alumbrado general & 132 & 1.00 & 132 & 0.67 & Incluye sala, dormitorios, bano, lavanderia y circulaciones \\
C2 & Tomacorrientes generales & 2520 & 0.70 & 1764 & 8.91 & Incluye sala-comedor, dormitorios, bano y lavanderia \\
C3 & Cocina & 1200 & 0.80 & 960 & 4.85 & Tomacorrientes para pequenos artefactos y refrigeracion \\
C4 & Lavanderia & 360 & 0.80 & 288 & 1.45 & Tomacorrientes y apoyo de servicio \\
C5 & Carga especial & 750 & 1.00 & 750 & 3.79 & Bomba de agua o carga de servicio equivalente \\
\textbf{Total} &  & \textbf{4962} &  & \textbf{3894} & \textbf{19.67} &  \\
\bottomrule
\end{tabularx}
\end{table}

\subsection{Levantamiento de cargas por ambiente}

\small
\setlength{\tabcolsep}{2pt}
\begin{longtable}{c p{2.6cm} p{2.8cm} p{2cm} c r r p{1.2cm} p{3cm}}
\caption{Levantamiento de cargas por ambiente}\label{tab:levantamiento-cargas}\\
\toprule
\textbf{Item} & \textbf{Ambiente} & \textbf{Equipo o punto} & \textbf{Tipo} & \textbf{Cant.} & \textbf{W unit.} & \textbf{W total} & \textbf{Circuito} & \textbf{Observacion} \\
\midrule
\endfirsthead
\toprule
\textbf{Item} & \textbf{Ambiente} & \textbf{Equipo o punto} & \textbf{Tipo} & \textbf{Cant.} & \textbf{W unit.} & \textbf{W total} & \textbf{Circuito} & \textbf{Observacion} \\
\midrule
\endhead
1 & Sala-comedor & Luminaria LED & Alumbrado & 2 & 12 & 24 & C1 & Iluminacion general del area social \\
2 & Sala-comedor & Tomacorriente doble & Tomacorrientes & 4 & 180 & 720 & C2 & Uso general y equipos domesticos ligeros \\
3 & Cocina & Luminaria LED & Alumbrado & 2 & 12 & 24 & C1 & Apoyo visual en ambiente de trabajo \\
4 & Cocina & Tomacorriente de pequeno artefacto & Tomacorrientes & 4 & 300 & 1200 & C3 & Refrigeradora, licuadora, microondas y similares \\
5 & Dormitorio principal & Luminaria LED & Alumbrado & 1 & 12 & 12 & C1 & Iluminacion puntual \\
6 & Dormitorio principal & Tomacorriente doble & Tomacorrientes & 3 & 180 & 540 & C2 & Carga de uso comun \\
7 & Dormitorio secundario 1 & Luminaria LED & Alumbrado & 1 & 12 & 12 & C1 & Iluminacion puntual \\
8 & Dormitorio secundario 1 & Tomacorriente doble & Tomacorrientes & 2 & 180 & 360 & C2 & Carga de uso comun \\
9 & Dormitorio secundario 2 & Luminaria LED & Alumbrado & 1 & 12 & 12 & C1 & Iluminacion puntual \\
10 & Dormitorio secundario 2 & Tomacorriente doble & Tomacorrientes & 2 & 180 & 360 & C2 & Carga de uso comun \\
11 & Bano & Luminaria LED & Alumbrado & 1 & 12 & 12 & C1 & Instalar fuera de zonas de contacto directo con agua \\
12 & Bano & Tomacorriente protegido & Tomacorrientes & 1 & 180 & 180 & C2 & Debe contar con proteccion y ubicacion adecuada \\
13 & Lavanderia & Luminaria LED & Alumbrado & 1 & 12 & 12 & C1 & Area de servicio \\
14 & Lavanderia & Tomacorriente doble & Tomacorrientes & 2 & 180 & 360 & C4 & Lavadora y apoyo de planchado o servicio \\
15 & Area tecnica & Bomba de agua o carga especial & Especial & 1 & 750 & 750 & C5 & Circuito dedicado \\
16 & Circulaciones & Luminaria LED & Alumbrado & 2 & 12 & 24 & C1 & Escalera y pasadizos \\
\bottomrule
\end{longtable}
\normalsize

\subsection{Interpretacion de la demanda}

La demanda total estimada de 3894 W equivale aproximadamente a 3.89 kW. Con una tension de 220 V y factor de potencia de 0.90, la corriente demandada es cercana a 19.67 A. En una propuesta academica para vivienda, este valor justifica un interruptor general de 32 A o 40 A, segun el criterio de reserva adoptado y la longitud del alimentador.

\section{Diseno del tablero electrico general y tableros de distribucion}

\subsection{Criterio general de sectorizacion}

Dada la existencia de tres niveles, la instalacion puede resolverse con un tablero general y tableros de distribucion por nivel funcional, con el fin de reducir recorridos, organizar los circuitos y facilitar mantenimiento. Para fines academicos se propone la siguiente estructura:

\begin{table}[H]
\centering
\caption{Propuesta de tableros de distribucion}
\label{tab:tableros-distribucion}
\begin{tabularx}{\textwidth}{p{2.2cm}p{4.8cm}p{4.5cm}X}
\toprule
\textbf{Elemento} & \textbf{Ubicacion propuesta} & \textbf{Funcion} & \textbf{Comentario tecnico} \\
\midrule
TG-01 & Primer piso, cerca del ingreso principal & Recibir el suministro, cortar la energia general y alimentar los tableros secundarios & Debe permanecer accesible y señalizado \\
TD-01 & Zona central de la vivienda, asociada a los niveles habitables & Distribuir circuitos de alumbrado y tomacorrientes generales & Sectoriza las cargas de uso diario \\
TD-02 & Tercer piso o area de servicio & Distribuir lavanderia y carga especial & Reduce longitud de circuito y mejora ordenamiento \\
\bottomrule
\end{tabularx}
\end{table}

\subsection{Propuesta de componentes del tablero general}

\begin{table}[H]
\centering
\caption{Componentes basicos del tablero general}
\label{tab:componentes-tablero}
\begin{tabularx}{\textwidth}{p{4cm}X}
\toprule
\textbf{Componente} & \textbf{Especificacion academica} \\
\midrule
Gabinete & Empotrado o superficial, grado de proteccion apto para interior residencial \\
Interruptor general & 2 polos, 40 A, curva residencial, maniobra manual \\
Interruptor diferencial & 2 polos, 40 A, 30 mA, tipo para proteccion de personas \\
Proteccion contra sobretensiones & Opcional como mejora academica, tipo 2 si el docente la solicita \\
Barra de neutro & Aislada y rotulada \\
Barra de tierra & Con conexion franca al sistema de puesta a tierra \\
Reserva & Espacios libres para ampliacion futura \\
\bottomrule
\end{tabularx}
\end{table}

\subsection{Distribucion funcional de los tableros}

\begin{table}[H]
\centering
\caption{Distribucion funcional de los tableros}
\label{tab:funcion-tableros}
\begin{tabularx}{\textwidth}{p{2cm}p{4cm}p{3cm}X}
\toprule
\textbf{Tablero} & \textbf{Circuitos alojados} & \textbf{Nivel de carga} & \textbf{Observacion} \\
\midrule
TG-01 & Alimentacion general y derivacion a TD-01 y TD-02 & Bajo a medio & Solo seccionamiento y proteccion principal \\
TD-01 & C1, C2 y C3 & Medio & Atiende alumbrado, tomacorrientes generales y cocina \\
TD-02 & C4 y C5 & Bajo a medio & Atiende lavanderia y carga especial \\
\bottomrule
\end{tabularx}
\end{table}

\subsection{Criterio de dimensionamiento del tablero}

Para el nivel de carga estimado, un tablero principal de 8 a 12 modulos utiles resulta suficiente, siempre que se reserve espacio para mantenimiento o ampliacion. En caso de emplear dos tableros de distribucion, cada uno debe contar con barra de neutro, barra de tierra y rotulacion individual de circuitos.

\section{Distribucion de circuitos}

\subsection{Circuitos propuestos}

\begin{table}[H]
\centering
\caption{Circuitos propuestos}
\label{tab:circuitos-propuestos}
\begin{tabularx}{\textwidth}{p{1.5cm}p{3.2cm}p{4cm}p{4cm}X}
\toprule
\textbf{Circuito} & \textbf{Uso} & \textbf{Ambientes atendidos} & \textbf{Proteccion preliminar} & \textbf{Conductores preliminares} \\
\midrule
C1 & Alumbrado general & Sala-comedor, cocina, dormitorios, bano, lavanderia y circulaciones & 10 A, 2 polos & 3 x 1.5 mm2 Cu \\
C2 & Tomacorrientes generales & Sala-comedor, dormitorios y bano & 16 A, 2 polos & 3 x 2.5 mm2 Cu \\
C3 & Cocina & Tomacorrientes de pequenos artefactos y refrigeracion & 20 A, 2 polos & 3 x 2.5 mm2 Cu o 4 mm2 Cu si la longitud lo exige \\
C4 & Lavanderia & Lavadora y tomacorrientes de servicio & 16 A, 2 polos & 3 x 2.5 mm2 Cu \\
C5 & Carga especial & Bomba de agua o equipo de servicio equivalente & 16 A, 2 polos & 3 x 2.5 mm2 Cu \\
\bottomrule
\end{tabularx}
\end{table}

\subsection{Criterios de distribucion por ambiente}

\begin{itemize}
  \item La sala-comedor concentra los puntos de uso social y varios tomacorrientes generales.
  \item La cocina se separa por su mayor exigencia de carga y por la presencia de equipos de uso discontinuo.
  \item Los dormitorios comparten un circuito de tomacorrientes generales por similitud de uso.
  \item El bano debe incorporar tomacorriente protegido y ubicacion fuera de zonas de salpicadura directa.
  \item La lavanderia se separa por la presencia de equipos con uso especial.
  \item La carga especial se reserva en circuito independiente para garantizar maniobra y proteccion selectiva.
\end{itemize}

\subsection{Observaciones de seguridad}

\begin{itemize}
  \item Los circuitos de tomacorrientes deben contar con conductor de proteccion.
  \item En ambientes humedos se recomienda tomacorriente con tapa o grado de proteccion adecuado.
  \item Los puntos de luz deben disponer de caja, derivacion accesible y conexion firme.
  \item Las cargas especiales no deben compartirse con tomacorrientes generales.
\end{itemize}

\section{Seleccion de conductores y dispositivos de proteccion}

\subsection{Criterio de seleccion}

La seleccion preliminar de conductores se realiza por capacidad de corriente, uso residencial tipico y reserva de seguridad. Debe verificarse, en un desarrollo posterior, la caida de tension, la temperatura ambiente, la forma de instalacion y la agrupacion de conductores.

\subsection{Seleccion preliminar de conductores}

\begin{table}[H]
\centering
\caption{Seleccion preliminar de conductores}
\label{tab:seleccion-conductores}
\begin{tabularx}{\textwidth}{p{4.2cm}p{3.4cm}p{2.6cm}X}
\toprule
\textbf{Tramo o circuito} & \textbf{Seccion propuesta} & \textbf{Material} & \textbf{Uso tecnico} \\
\midrule
Alimentacion general desde acometida a tablero general & 2 x 10 mm2 Cu + PE & Cobre & Reserva adecuada para corriente general y crecimiento futuro \\
Derivacion a TD-01 & 2 x 6 mm2 Cu + PE & Cobre & Alimentacion de circuito sectorizado \\
Derivacion a TD-02 & 2 x 6 mm2 Cu + PE & Cobre & Alimentacion de circuito sectorizado \\
C1 alumbrado & 3 x 1.5 mm2 Cu & Cobre & Circuito de iluminacion \\
C2 tomacorrientes generales & 3 x 2.5 mm2 Cu & Cobre & Cargas domesticas generales \\
C3 cocina & 3 x 2.5 mm2 Cu & Cobre & Cargas de pequeno artefacto \\
C4 lavanderia & 3 x 2.5 mm2 Cu & Cobre & Lavadora y servicio \\
C5 carga especial & 3 x 2.5 mm2 Cu & Cobre & Bomba de agua o equipo equivalente \\
Puesta a tierra & 6 mm2 Cu minimo, sujeto a criterio final & Cobre & Conexion principal entre barra de tierra y electrodo \\
\bottomrule
\end{tabularx}
\end{table}

\subsection{Protecciones preliminares}

\begin{table}[H]
\centering
\caption{Protecciones preliminares}
\label{tab:protecciones-preliminares}
\begin{tabularx}{\textwidth}{p{4cm}p{5cm}X}
\toprule
\textbf{Circuito} & \textbf{Interruptor termomagnetico propuesto} & \textbf{Observacion} \\
\midrule
Alimentacion general & 2P, 40 A & Proteccion y seccionamiento principal \\
TD-01 & 2P, 25 A & Proteccion del subtablero \\
TD-02 & 2P, 25 A & Proteccion del subtablero \\
C1 alumbrado & 2P, 10 A & Proteccion por baja carga y conductores de 1.5 mm2 \\
C2 tomacorrientes generales & 2P, 16 A & Proteccion residencial tipica \\
C3 cocina & 2P, 20 A & Mayor demanda por pequenos artefactos \\
C4 lavanderia & 2P, 16 A & Equipo de uso continuo parcial \\
C5 carga especial & 2P, 16 A & Circuito dedicado \\
\bottomrule
\end{tabularx}
\end{table}

\subsection{Coordinacion conductor-proteccion}

La relacion basica entre corriente de diseno, corriente nominal del interruptor y capacidad del conductor debe mantener el criterio:

\begin{equation}
I_b \leq I_n \leq I_z
\end{equation}

Donde:

\begin{itemize}
  \item $I_b$: corriente de diseno del circuito.
  \item $I_n$: corriente nominal de la proteccion.
  \item $I_z$: capacidad admisible del conductor.
\end{itemize}

Con esta relacion se evita que el dispositivo de proteccion permita una sobrecarga superior a la capacidad del conductor, reduciendo el riesgo de sobretemperatura.

\subsection{Observaciones tecnicas adicionales}

\begin{itemize}
  \item En circuitos de cocina y lavanderia se recomienda prever reservas de capacidad por el uso de electrodomesticos portatiles.
  \item Si la longitud de algun circuito supera de manera significativa el recorrido previsto, la seccion debera incrementarse para controlar la caida de tension.
  \item La seleccion final debe contrastarse con tablas del fabricante y con el criterio del CNE-U para instalacion en tuberia o canalizacion equivalente.
\end{itemize}

\section{Sistema de puesta a tierra}

\subsection{Funcion del sistema}

El sistema de puesta a tierra tiene por finalidad establecer una referencia de potencial y facilitar la derivacion de corrientes de falla, reduciendo el riesgo de choque electrico por contacto indirecto. En viviendas, este sistema debe estar integrado a tomacorrientes, carcasas metalicas y barras de tierra del tablero.

\subsection{Componentes minimos}

\begin{table}[H]
\centering
\caption{Componentes minimos del sistema de puesta a tierra}
\label{tab:componentes-puesta-tierra}
\begin{tabularx}{\textwidth}{p{4cm}X}
\toprule
\textbf{Componente} & \textbf{Especificacion academica} \\
\midrule
Electrodo de puesta a tierra & Varilla copperweld o equivalente, instalada en zona accesible para inspeccion \\
Caja de registro & Permite inspeccion y mantenimiento del electrodo \\
Conductor principal de tierra & Cobre de seccion adecuada, conectado a la barra de tierra del tablero \\
Barra de tierra & Barra independiente, rotulada y unida a los conductores de proteccion \\
Conductores de proteccion & Conectan tomacorrientes y partes metalicas accesibles con la barra de tierra \\
Conectores & Abrazaderas o conectores certificados para uniones firmes y durables \\
\bottomrule
\end{tabularx}
\end{table}

\subsection{Criterio de dimensionamiento}

Para una propuesta academica residencial se adopta como criterio preliminar:

\begin{itemize}
  \item Conductor principal de puesta a tierra: 6 mm2 Cu minimo.
  \item Conductores de proteccion de circuitos derivados: 2.5 mm2 Cu como referencia minima para tomacorrientes.
  \item Continuidad metalica entre barra de tierra, electrodo y equipos conectados.
  \item Medicion de resistencia de puesta a tierra en una fase posterior del proyecto, si se ejecutara.
\end{itemize}

\subsection{Recomendaciones de instalacion}

\begin{itemize}
  \item El electrodo debe quedar protegido contra deterioro mecanico y corrosivo.
  \item La barra de tierra debe ubicarse dentro del tablero o en gabinete adyacente.
  \item La puesta a tierra no debe compartirse con conductores de neutro en el interior del tablero, salvo la configuracion especifica del sistema y el criterio normativo aplicable.
  \item Toda conexion debe quedar firme, accesible y rotulada.
\end{itemize}
